% ============================================================================
% WILSON Fire Project — Validation & Benchmark Data Appendix
% Generated 2026-07-08 from automated test runs (scripts + raw JSON in the
% session scratchpad; hindcast runs use FIRED perimeters + ERA5 weather).
% Compile standalone, or \input individual tables into the main report.
% ============================================================================
\documentclass[11pt]{article}
\usepackage[margin=2.2cm]{geometry}
\usepackage{booktabs}
\usepackage{caption}
\usepackage{amsmath}
\begin{document}

\section*{WILSON Fire Project --- Validation \& Benchmark Data}

\subsection*{A. Hindcast validation against FIRED satellite perimeters}

Thirteen end-to-end hindcast runs (NASA FIRED daily perimeters as ground
truth, Open-Meteo ERA5 weather, Copernicus COP30 terrain + CORINE fuels).
Differential Evolution calibrates five hidden parameters (moisture offset,
wind multiplier, wind direction offset, ignition threshold, burn-time
multiplier); budget \texttt{maxiter}=2, \texttt{popsize}=4; coarse grid
$200{\times}200$ (${\sim}278$\,m), final IoU scored on $800{\times}800$
(${\sim}69$\,m). \emph{Point} = seeded from day-1 polygon centroids
(from-ignition skill). \emph{Assim.}~$d_n$ = seeded from the observed
cumulative perimeter of day $n$ (FARSITE-style validation), predicting the
next 24\,h (48\,h where noted).

\begin{table}[htbp]
\centering
\caption{Hindcast IoU for all scenarios. $A_p/A_t$ = predicted/truth area
ratio; $k_w$ = calibrated wind multiplier; $\theta_i$ = ignition threshold;
$b_t$ = burn-time multiplier.}
\small
\begin{tabular}{llrrrrrl}
\toprule
Fire & Scenario & IoU (\%) & $A_p/A_t$ & $k_w$ & $\theta_i$ & $b_t$ &
Interpretation \\
\midrule
Rhodes 2023 & Point, $d_1$      & 42.6 & 1.9 & 0.62 & 0.55 & 2.57 & $d_1$ already 968\,ha, spreading \\
Rhodes 2023 & Assim.\ $d_2$     & 20.8 & 4.6 & 1.04 & 0.96 & 1.75 & real $d_2{\to}d_3$ growth tiny (+17\%) \\
Rhodes 2023 & Assim.\ $d_4$     & 47.6 & 1.8 & 0.74 & 0.80 & 1.78 & main run of fire \\
Evoia 2021  & Point, $d_1$      &  1.9 & 26  & 0.76 & 0.61 & 1.03 & $d_1$ smolder (43\,ha), unmatchable \\
Evoia 2021  & Assim.\ $d_2$     &  9.4 & 9.6 & 2.20 & 0.16 & 1.66 & $d_2{\to}d_3$ eruption ($43\times$) \\
Evoia 2021  & Assim.\ $d_4$     & 53.2 & 1.1 & 1.41 & 0.91 & 3.46 & eruption phase, well captured \\
Evoia 2021  & Assim.\ $d_4$, 48\,h & 52.9 & 1.6 & 2.53 & 0.72 & 4.14 & 2-day forecast holds up \\
Evoia 2021  & Assim.\ $d_6$     & 66.7 & 1.4 & 1.13 & 0.88 & 1.03 & steadier growth phase \\
Evoia 2021  & Assim.\ $d_8$     & 90.2 & 1.1 & 1.10 & 0.76 & 1.46 & late stage, small marginal growth \\
Evros 2023  & Point, $d_1$      &  2.0 & 24  & 2.49 & 0.59 & 1.18 & $d_1$ smolder (22\,ha) \\
Evros 2023  & Assim.\ $d_4$     & 43.0 & 2.2 & 1.59 & 0.47 & 4.55 & explosive $d_4{\to}d_5$ ($3\times$) \\
Evros 2023  & Assim.\ $d_8$     & 71.3 & 1.4 & 1.04 & 0.96 & 1.75 & mid--late stage \\
\midrule
\multicolumn{2}{l}{Mean, assimilated $d_4$ (3 fires)} & 47.9 & & & & & \\
\multicolumn{2}{l}{Mean, all assimilated (9 runs)}    & 50.6 & & & & & \\
\multicolumn{2}{l}{Mean, point $d_1$ (3 fires)}       & 15.5 & & & & & \\
\bottomrule
\end{tabular}
\end{table}

\noindent\textbf{Reading the sweep.} Skill tracks the \emph{relative}
next-day growth of the real fire: seeding during eruption days
(Evoia $d_2{\to}d_3$: $43\times$ area jump) or stall days (Rhodes
$d_2{\to}d_3$: $+17\%$, while the unsuppressed model keeps running) hurts;
mid/late-stage seeding reaches 43--90\%. The model \emph{always}
over-predicts ($A_p/A_t>1$ in every run): no suppression or evacuation
effects are modelled, and a low DE budget was used. A repeat of Evoia
$d_4$ at half the DE budget scored 51.6\% vs.\ 53.2\%, so results are not
budget-limited. Note the cumulative-scoring convention includes the
assimilated seed area, which inflates late-day IoU.

\subsection*{B. Historical fire context (FIRED + ERA5)}

\begin{table}[htbp]
\centering
\caption{FIRED cumulative burned area (ha) by day, and ERA5 daily-mean
weather at the ignition-day / day-4 seeds. Low ERA5 wind on eruption days
(0.4--3\,m/s) is why the optimizer's wind multiplier and the $6{\times}$
upper bound exist: the 0.25$^\circ$ reanalysis misses mesoscale and
pyroconvective winds.}
\begin{tabular}{lrrrrrrr}
\toprule
 & \multicolumn{4}{c}{Cumulative area (ha)} & \multicolumn{3}{c}{ERA5 (day 1 / day 4)} \\
\cmidrule(lr){2-5}\cmidrule(lr){6-8}
Fire & $d_1$ & $d_2$ & $d_4$ & final & wind (m/s) & $T$ ($^\circ$C) & RH (\%) \\
\midrule
Evoia 2021 (ev.\ 2871, 12\,d)  & 43   & 130   & 11\,226 & 52\,521 & 1.5 / 0.4 & 27.9 / 34.0 & 61 / 22 \\
Rhodes 2023 (ev.\ 130111, 10\,d) & 968 & 1\,763 & 4\,174 & 19\,874 & 1.7 / 3.0 & 26.9 / 24.9 & 52 / 76 \\
Evros 2023 (ev.\ 105620, 18\,d) & 22   & 1\,789 & 17\,928 & 89\,422 & 1.5 / 2.5 & 22.7 / 24.3 & 85 / 55 \\
\bottomrule
\end{tabular}
\end{table}

\begin{table}[htbp]
\centering
\caption{Fuel composition of each $0.5^\circ{\times}0.5^\circ$ simulation
domain (CORINE + OSM overlay, $200{\times}200$ grid; non-combustible
includes sea for these coastal domains).}
\begin{tabular}{lrrrrrr}
\toprule
Fire & Non-comb. & Roads/urban & Pine & Maquis & Grass/phrygana & Oak \\
\midrule
Evoia 2021  & 45.9\% & 16.5\% & 11.0\% & 6.7\% & 4.6\% & 2.4\% \\
Rhodes 2023 & 58.5\% & 12.0\% &  6.4\% & 9.0\% & 8.6\% & --   \\
Evros 2023  & 47.7\% & 13.5\% &  6.1\% & 6.4\% & 10.9\% & 9.8\% \\
\bottomrule
\end{tabular}
\end{table}

\subsection*{C. Fire-model physics characterization}

Controlled synthetic landscapes (uniform fuel, flat unless noted, 50\,m
cells, $dt=0.5$\,min, 60-min fires, fixed RNG seed).

\begin{table}[htbp]
\centering
\caption{Rate of spread (dry grass, moisture 0.06). Left: vs.\ open wind
(flat). Centre: vs.\ slope (no wind). Right: vs.\ moisture (6\,m/s wind).
Downwind ROS saturates above ${\sim}12$\,m/s open wind (Rothermel
1\,500\,ft/min effective-wind cap); upwind spread stays at the zero-wind
0.8\,m/min backing rate. Moisture shows the $\eta_m$ extinction cliff
between 0.10 and 0.15 (grass $m_x$).}
\begin{tabular}{rr@{\qquad}rr@{\qquad}rr}
\toprule
Wind (m/s) & ROS (m/min) & Slope ($^\circ$) & ROS$_{\uparrow}$ (m/min) & Moisture & ROS (m/min) \\
\midrule
 0 &  0.8 &  0 &  0.8 & 0.03 & 39.2 \\
 2 &  8.3 &  5 &  0.8 & 0.06 & 31.7 \\
 4 & 18.3 & 10 &  3.3 & 0.10 & 28.3 \\
 6 & 31.7 & 15 &  5.0 & 0.15 &  0.8 \\
 9 & 64.2 & 20 &  7.5 & 0.20 &  0.8 \\
12 & 82.5 & 30 & 15.0 & 0.24 &  0.8 \\
15 & 82.5 &    &      &      &      \\
\bottomrule
\end{tabular}
\end{table}

\begin{table}[htbp]
\centering
\caption{Fuel-dependent behaviour at 6 and 12\,m/s open wind. The wind
adjustment factor (WAF) reduces open wind to midflame wind under canopy;
surface spread under pine/oak canopy is therefore minimal at ERA5-scale
winds --- crown-fire behaviour is absorbed by the optimizer's wind and
per-fuel multipliers instead.}
\begin{tabular}{lrrrr}
\toprule
Fuel & WAF & Midflame @6\,m/s & ROS @6 (m/min) & ROS @12 (m/min) \\
\midrule
Dry grass          & 0.40 & 2.40\,m/s & 31.7 & 82.5 \\
Maquis dense shrub & 0.20 & 1.20\,m/s & 18.3 & 45.0 \\
Phrygana low scrub & --   & --        & 17.5 & --   \\
Olive grove        & --   & --        &  6.7 & --   \\
Aleppo pine        & 0.10 & 0.60\,m/s &  0.8 &  1.7 \\
Oak forest         & 0.12 & 0.72\,m/s &  0.8 &  3.3 \\
\bottomrule
\end{tabular}
\end{table}

\begin{table}[htbp]
\centering
\caption{Numerical properties. Resolution test: identical
$10{\times}10$\,km domain, 2-h fire, 6\,m/s wind, CFL-scaled
$dt \propto \Delta x$. Areas agree within 3\% between 100- and 50-m cells;
the 25-m grid over-shoots ${\sim}40\%$ because firebrand spotting
distances are resolution-independent in metres but land more, finer
ignition targets. Determinism: two runs with the same NumPy seed produce
bit-identical states.}
\begin{tabular}{lrrr@{\qquad}lr}
\toprule
Grid & Cell (m) & $dt$ (min) & Burned (ha) & Performance & ms/step \\
\midrule
$100^2$ & 100 & 1.00 & 677 & $200^2$ grid & 0.9 \\
$200^2$ &  50 & 0.50 & 696 & $400^2$ grid & 3.5 \\
$400^2$ &  25 & 0.25 & 992 & $800^2$ grid & 15.8 \\
\bottomrule
\end{tabular}
\end{table}

\noindent At the interactive default ($800^2$, $dt=0.05$\,min) the engine
sustains ${\sim}63$ steps/s $\approx$ 3.2 simulated minutes per wall second;
the $200^2$ optimizer grid runs ${\sim}1\,200$ steps/s. WebSocket frame
encoding (RLE GeoJSON, 80\,k fire cells on $800^2$) costs 7.3\,ms and
104\,KB per frame. Microclimate memory serializes to 1.4\,MB JSON in
25\,ms with a lossless round-trip.

\subsection*{D. Intervention efficacy}

Identical wind-driven grass fires (6\,m/s, $300^2$ grid); barrier placed
3\,km downwind of ignition; burned area after 3\,h vs.\ control.

\begin{table}[htbp]
\centering
\caption{Burned area after 3\,h. The ember-transport experiment (bottom)
isolates why containment lines underperform: with spotting disabled the
0.95-strength line holds completely, so \emph{all} line breaching is
firebrand transport over the corridor --- consistent with real
ember-driven Mediterranean fires. Firebreaks block embers whose flight
path crosses them; water drops on the head knock the fire back directly.}
\begin{tabular}{lrr}
\toprule
Scenario & Burned (ha) & Reduction \\
\midrule
Control (no intervention)               & 1\,581 & --    \\
Firebreak, 150\,m corridor              &   750  & 53\%  \\
Firebreak, 250\,m corridor              &   734  & 54\%  \\
Containment line, $s=0.7$               & 1\,493 &  6\%  \\
Containment line, $s=0.95$              & 1\,404 & 11\%  \\
Containment $s=0.7$ + water + humidity  & 1\,470 &  7\%  \\
3 aerial drops on fire head at 60\,min  &   861  & 46\%  \\
\midrule
Containment $s=0.95$, spotting disabled &   743  & 53\% (line holds; fire never crosses) \\
\bottomrule
\end{tabular}
\end{table}

\subsection*{E. Sandbox PSO optimizer}

Particle-swarm calibration (12 particles $\times$ 8 iterations,
12-parameter space, synthetic $200^2$ terrain, 2-h fire vs.\ a drawn truth
polygon; 6 worker threads).

\begin{table}[htbp]
\centering
\caption{Best-so-far IoU per PSO iteration. The microclimate warm start
(ignition-frequency prior from 50 previous optimizer simulations)
accelerates early convergence but converges to a worse optimum here: the
learned spatial prior biases the search toward previously-burned corridors.
Useful for fast triage; disable for final calibration.}
\begin{tabular}{lrrrrrrrr@{\qquad}rr}
\toprule
 & \multicolumn{8}{c}{Iteration} & Final & Wall (s) \\
\cmidrule(lr){2-9}
Run & 1 & 2 & 3 & 4 & 5 & 6 & 7 & 8 & IoU & \\
\midrule
Fresh                  & .00 & .37 & .40 & .40 & .44 & .45 & .50 & .52 & \textbf{0.556} & 149 \\
Microclimate warm start& .00 & .42 & .46 & .47 & .47 & .47 & .47 & .47 & 0.476 & 129 \\
\bottomrule
\end{tabular}
\end{table}

\subsection*{F. Software verification \& project statistics}

\begin{itemize}
  \item Regression suite (optimizer state-freeze, canopy-factor
        persistence, cross-resolution fire restore, parameter
        application): \textbf{5/5 pass}.
  \item All Python modules byte-compile; WebSocket protocol handlers
        import cleanly.
  \item Determinism: bit-identical simulation states under a fixed RNG
        seed (Table in \S C).
  \item Codebase: ${\sim}9\,100$ lines of Python (core model, pipelines,
        server) + 5\,300-line web client; 23 Greek fuel classes (4
        non-burning); 41 generated UI icons.
  \item Data assets: FIRED Greece 2000--2024 daily perimeters (22\,MB
        GeoPackage, 24\,867 features); 34 cached DEM/CORINE rasters.
\end{itemize}

\subsection*{G. Defects found and fixed during this test campaign}

\begin{enumerate}
  \item Firebrand spotting ignored containment-line strength at the
        landing cell; brands now damped by local crew strength (embers
        beyond the corridor unaffected --- Table \S D reflects the fix).
  \item Hindcast: ERA5 weather was fetched for the FIRMS query start date
        rather than the seed day whenever \texttt{ignition\_day} $>1$.
  \item Hindcast: the perimeter-assimilation + tunable-persistence
        upgrade existed only in an unpushed working copy; re-implemented
        and validated (Evoia $d_4$: 53.2\% vs.\ 52.0\% previously).
  \item Cosmetic: results report printed physically impossible negative
        ``effective midflame moisture'' (clip was not applied in the
        printout).
\end{enumerate}

\subsection*{H. Overall assessment}

From-ignition (day-1) forecasting is only skilful when the fire is
already large at first detection (Rhodes: 42.6\%); smoldering-start fires
are structurally over-predicted (${\sim}2\%$). Operational-style
assimilated forecasting averages ${\sim}48\%$ IoU at day 4 and 71--90\%
late-event, inside the 30--60\% band typical of published FARSITE/FlamMap
validations, with a consistent over-prediction bias attributable to
missing suppression modelling and smooth ERA5 winds. The physics engine
reproduces the qualitative Rothermel responses (wind power-law with
saturation, slope amplification, moisture extinction cliff, canopy wind
sheltering), is deterministic under a fixed seed, resolution-consistent
to within 3\% down to 50-m cells, and fast enough for interactive use.

\end{document}
